% Source: https://artofproblemsolving.com/community/c3h1681395_circles_and_incentres
% Stable identity: aops-topic-1681395-q01
Let $O$ be the centre. The whole figure is symmetric in line $AO$, so the midpoint $D$ of the minor arc $BC$ lies on $AO$, and $AD$ is the bisector of $\angle BAC$. The minor arc $BC$ lies inside triangle $ABC$, so $D$ is inside the triangle.

\begin{asy}
usepackage("amsmath");
size(4.5cm);
defaultpen(fontsize(10pt));
pair O=(0,0);real R=1;pair A=(0,-2.4);
real t=aSin(R/2.4);
pair B=R*dir(-90+ (90-t)),C=R*dir(-90-(90-t));
pair D=(0,-R);
draw(circle(O,R),gray);
draw(A--B--C--cycle,linewidth(0.9));
draw(A--D,dashed+blue);draw(B--D,blue);draw(C--D,blue);
dot(O);dot(D);
label("$O$",O,N);label("$A$",A,S);label("$B$",B,E);label("$C$",C,W);label("$D$",D,dir(-45));
label("\mbox{}",(min(currentpicture,user=true)+max(currentpicture,user=true))/2,black); // keeps the text after the figure black when LaTeX typesets the labels inline
\end{asy}

By the tangent-chord angle theorem, the angle between tangent $BA$ and chord $BD$ equals the inscribed angle on the arc $BD$ from the other side:
\[
\angle ABD=\angle BCD .
\]
Since arcs $BD$ and $DC$ are equal, the inscribed angles on them are equal: $\angle BCD=\angle DBC$. Hence $\angle ABD=\angle DBC$, so $BD$ bisects $\angle ABC$.

$D$ lies on the bisectors of $\angle A$ and $\angle B$, so $D$ is the incentre of triangle $ABC$. $\blacksquare$
